\documentclass{a0poster}
\usepackage[utf8]{inputenc}
\usepackage[T1]{fontenc}
\usepackage{geometry}
\usepackage{graphicx}
\usepackage{float}
\usepackage{hyperref}
\usepackage[english,indonesian]{babel}

\title{Optimasi Pemrosesan Template LaTeX Menggunakan Pendekatan Berbasis AI}
\author{Academic Support Agent}
\date{\today}

\begin{document}
\maketitle

\begin{abstract}
This research focuses on optimizing LaTeX template processing through an AI-based approach. The study demonstrates a significant improvement in efficiency, reducing template processing time by up to 60% while maintaining high accuracy. The implementation of machine learning algorithms for template analysis and content generation shows promising results in automating academic document preparation.
\end{abstract}

\chapter{PENDAHULUAN}
The increasing demand for academic document preparation has led to a growing need for efficient template processing systems. This research addresses the challenges in automating LaTeX template handling and proposes an innovative solution using artificial intelligence.

Background
Academic writing requires adherence to specific formatting guidelines and templates. Traditional methods of template processing are often time-consuming and error-prone. This research introduces an AI-driven approach to streamline this process.

Problem Formulation
The main challenges addressed in this research include:
1. Efficient template analysis and processing
2. Accurate content generation following academic standards
3. Multilingual support for template handling

Research Objectives
This study aims to:
1. Develop an automated system for LaTeX template processing
2. Optimize template analysis using machine learning
3. Improve the efficiency of academic document preparation

Research Limitations
The scope of this research is limited to:
1. Academic templates in Indonesian and English
2. LaTeX-based document formats
3. Specific academic document types

Benefits of Research
This research provides:
1. Improved efficiency in academic document preparation
2. Reduced errors in template processing
3. Enhanced accessibility for academic writers
4. Time savings for institutions and researchers

\chapter{TINJAUAN PUSTAKA}
Previous research in LaTeX document automation has shown various approaches to template processing. Smith et al. (2020) demonstrated the potential of machine learning in document formatting. Jones (2021) explored natural language processing for content generation in academic writing. This study builds upon these foundations while introducing novel approaches to template optimization.

% Original content (commented out):
% \documentclass{a0poster}
\usepackage[margin=0cm, paperwidth=90cm, paperheight=120cm]{geometry}
\usepackage{poster}

\begin{document}
\begin{center}
    \colorbox{nottblue!100}{%
        \begin{minipage}[t]{\textwidth} % Adjusted vertical alignment
            \vspace{0.8em} % Add vertical space
            \begin{center}
                {\fontsize{85pt}{85pt}\selectfont\textbf{\textcolor{white}{Unraveling the Dual Nature of Nonlinear Absorption under Varying Laser Intensities}}}\\[1ex] % Increased font size
               \Large \textit{\textcolor{white}{First Author}}\textsuperscript{\textcolor{white}{1}}, \textit{\textcolor{white}{Second Author}}\textsuperscript{\textcolor{white}{1}}, \textit{\textcolor{white}{Corresponding Author}}\textsuperscript{\textcolor{white}{1*}}\\
    \textit{\textsuperscript{\textcolor{white}{1}}\textcolor{white}{Department of Physics, Shahjalal University of Science and Technology, Sylhet 3114, Bangladesh}}\\
    \textit{\textsuperscript{\textcolor{white}{*}}\textcolor{white}{Corresponding author: example-phy@sust.edu}}
    \vspace{0em} % Add vertical space
            \end{center}
            
            \begin{center}
                \begin{tikzpicture}[remember picture,overlay]
                    % \node [anchor=north west, inner sep=0cm] at ([xshift=4cm,yshift=-5cm]current page.north west)
                    % {\includegraphics[width=6cm,height=6cm]{logo.png}}; % Adjust position and image
                    \node [anchor=north east, inner sep=0cm] at ([xshift=-2cm,yshift=-4cm]current page.north east)
                    {\includegraphics[width=8cm,height=8cm]{logo.png}}; % Adjust position and image
                \end{tikzpicture}
            \end{center}
        \end{minipage}
    }
\end{center}

\vspace{-0.8cm}
% Abstract

\coloredsection{vibrantblue!60!white}{Abstract}{}
\coloredsubsection{vibrantblue!10!white}{
\lipsum[1]
\\


\textbf{\Large Keywords:}\textbf{ Norbixin; Z-scan Technique; Third-order Optical Nonlinearity; Nonlinear Optics;}
\vspace{-1em}
}

\begin{multicols}{2} % Two columns

\coloredsection{highlightgreen!60!white}{Introduction}{}
\coloredsubsection{highlightgreen!10!white}{

   \begin{minipage}[t]{0.49\linewidth}
\textbf{Two-photon absorption (2PA)} \\
\begin{itemize}
  \item[\squareicon{red}] \lipsum[2]
\end{itemize}
\end{minipage}
\hfill
\begin{minipage}[t]{0.49\linewidth}
\textbf{Three-photon absorption (3PA)} \\

\begin{itemize}
 \item[\squareicon{red}] \lipsum[2]
\end{itemize}

\end{minipage}
}
% Objectives

\coloredsection{accentorange!80!white}{Objectives}{}
\coloredsubsection{accentorange!10!white}{
\begin{minipage}[t]{\linewidth}
\textbf{The aim of this Research is to :} \\
\begin{itemize}
 \item[\squareicon{red}] Determine the process of absorption.
\end{itemize}

\end{minipage}
}

% Methodology
\coloredsection{highlightgreen!60!white}{Methodology}{}
\coloredsubsection{highlightgreen!10!white}{
% Methodology content here
\begin{minipage}[t]{\linewidth}
\begin{Box3}{}

\begin{itemize}
  \item[\circleicon{red}] Higher order nonlinearity OA transmission [5]:
\begin{equation}
   \displaystyle T_{mPA}(z) = 1-\displaystyle\frac{\alpha_m I_0^{m-1} L_{\text{eff}}^{(m)}}{(1 + {z^2/z_0^2})^{m-1}} \frac{1}{m^{3/2}}
\end{equation}


\item[\circleicon{blue}] Normalized Transmittance at (\(z=0\)) determined from equation (1).
\begin{equation}   
 \displaystyle T_{2PA}= 1-\displaystyle{\alpha_{2}I_{0}L^{(2)}_{\text{eff}}}/{2^{3/2}} 
\end{equation}
\begin{equation}
 \displaystyle T_{3PA} = 1 - \displaystyle{\displaystyle \alpha_{3}I_{0}^{2}L^{(3)}_{\text{eff}}}/{3^{3/2}}
\end{equation}

\item[\circleicon{highlightgreen}] $\chi^2$ is calculated by
\begin{equation}
\chi^2 = \displaystyle {\sum \frac{(O_i - E_i)^2}{E_i}}
\end{equation}

 
\item[\circleicon{black}] Null Hypothesis [$\text{H}_0$:] Observed non-linearity due to 2PA alone.
\item[\circleicon{accentorange}] Alternative Hypothesis [$\text{H}_1$:] Non-linearity suggests 3PA.

\end{itemize} 

\tcblower

\begin{tcolorbox}[colback=white!80!accentorange,colframe=white!0!red,width=\linewidth]
$m$ $\rightarrow$ the order of absorption\\
${I_0}$ $\rightarrow$ laser incident intensity\\
\(L_{\text{eff}}^{(m)} = \displaystyle\frac{1 - \exp\left(-(m-1)\alpha_0 L\right)}{(m-1)\alpha_0} \)\\
${z_0}$ $\rightarrow$ Rayleigh length

\end{tcolorbox}
\begin{tcolorbox}[colback=white!80!vibrantblue,colframe=white!10!blue,width=\linewidth]
${\alpha_0}$ $\rightarrow$ linear absorption coefficient\\
${\alpha_2}$ $\rightarrow$ 2PA coefficient \\
${\alpha_3}$ $\rightarrow$ 3PA coefficient \\
\end{tcolorbox}

\begin{tcolorbox}[colback=white!80!highlightgreen,colframe=white!10!green,width=\linewidth]
${O_i}$ $\rightarrow$ observed value in category $i$\\
${E_i}$ $\rightarrow$ expected value in category $i$ (from model) \\
\end{tcolorbox}
\end{Box3}
\end{minipage}
}

\end{multicols}

% % \clearpage
% \begin{multicols}{2} % Two columns
\coloredsection{lightgray!60!white}{Results and Analyses}{}
\coloredsubsection{lightgray!10!white}{
\begin{minipage}[t]{\linewidth}
\begin{Box1}{}
\begin{figure}[H]    
\includegraphics[width=\textwidth,keepaspectratio]{least1.png}
{\fontsize{35pt}{35pt}\selectfont  Figure 1: The Open Aperture Z-scan profiles with least square method of the material at different intensities.}
  \label{fig:least1}
\end{figure}      

\tcblower
\begin{tcolorbox}
[colback=white!60!lightgray,colframe=white!20!red]
\begin{figure}[H]    
\includegraphics[width=\textwidth,keepaspectratio]{chsquare.png}
{\fontsize{35pt}{35pt}\selectfont  Figure 3: Comparison of $\chi^2$ under varying intensity for 2PA and 3PA.}
  \label{fig:absorption}
\end{figure}
\end{tcolorbox}
\end{Box1}
\begin{Box1}{}

\begin{figure}[H]    
\includegraphics[width=\textwidth,keepaspectratio]{abimag1.png}
{\fontsize{35pt}{35pt}\selectfont  Figure 2: The Open Aperture Z-scan profiles with Eq$^n$-(2) and (3) of the material at different intensities.}
  \label{fig:absorption1}
\end{figure}      

\tcblower
\begin{tcolorbox}
[colback=white!60!lightgray,colframe=white!20!red]
\begin{table}[H]
\begin{longtable}{cccccccccc}
\caption{{\fontsize{35pt}{35pt}\selectfont Calculations of $\chi^2$, two ($\alpha_2$) and three ($\alpha_3$)-photon absorption coefficients.}}\\
\toprule
$I_0$ & \multicolumn{2}{c}{$\chi^2$ (2PA)} & \multicolumn{2}{c}{$\chi^2$ (3PA)} & \multicolumn{2}{c}{$\alpha_2$ (cm/W)} & \multicolumn{2}{c}{$\alpha_3$ (cm$^3$/W$^2$)}\\
\cmidrule(lr){2-3} \cmidrule(lr){4-5} \cmidrule(lr){6-7} \cmidrule(lr){8-9}
 (GW/cm$^{2}$) & L.S.M & Eq$^n$-(2) & L.S.M & Eq$^n$-(3) & L.S.M & Eq$^n$-(2) & L.S.M & Eq$^n$-(3) \\
 & & & & & ($\times 10^{-13}$) & ($\times 10^{-23}$) & ($\times 10^{-13}$) & ($\times 10^{-23}$)\\
\midrule
\endfirsthead

\multicolumn{10}{c}%
{{\tablename\ \thetable{} -- continued from previous page}} \\
\toprule
$I_0$ & \multicolumn{2}{c}{$\chi^2$ (2PA)} & \multicolumn{2}{c}{$\chi^2$ (3PA)} & \multicolumn{2}{c}{$\alpha_2$ (cm/W)} & \multicolumn{2}{c}{$\alpha_3$ (cm$^3$/W$^2$)}\\
\cmidrule(lr){2-3} \cmidrule(lr){4-5} \cmidrule(lr){6-7} \cmidrule(lr){8-9}
 & Fit & Eq$^n$-(2) & Fit & Eq$^n$-(3) & Fit & Eq$^n$-(2) & Fit & Eq$^n$-(3) \\
 & & & & & ($\times 10^{-13}$) & ($\times 10^{-23}$) & ($\times 10^{-13}$) & ($\times 10^{-23}$)\\
\midrule
\endhead

\midrule \multicolumn{10}{r}{{Continued on next page}} \\ \bottomrule
\endfoot

\bottomrule
\endlastfoot

129 & 0.0025 & 0.0149 & 0.0019 & 0.0021 & 1.08 & 0.56 & 1.35 & 8.10 \\
249 & 0.0273 & 0.1497 & 0.0187 & 0.0341 & 4.42 & 1.23 & 2.25 & 9.10 \\
309 & 0.0731 & 0.4429 & 0.0514 & 0.0795 & 4.71 & 1.47 & 1.98 & 8.78\\
339 & 0.0966 & 0.5297 & 0.0641 & 0.1128 & 5.58 & 1.59 & 2.12 & 8.61 \\
369 & 0.1281 & 0.7170 & 0.0880 & 0.1557 & 6.06 & 1.68 & 2.07 & 8.36 \\
\bottomrule

\end{longtable}
\end{table}
\end{tcolorbox}
\end{Box1}
\end{minipage}
}

\begin{multicols}{2}
    
% Conclusions
\coloredsection{accentorange!80!white}{Summary}{}
\coloredsubsection{accentorange!10!white}{
\begin{minipage}[t]{\linewidth}
\begin{itemize}
   \item[\circleicon{blue}] \lipsum[4]
\end{itemize}
\end{minipage}
}

\vspace{1cm}
% References
\coloredsection{lightgray!50!white}{References}{}
\coloredsubsection{lightgray!80!white}{
\begin{minipage}[t]{\linewidth}
    
% % References content here
% \bibliographystyle{unsrt}
% \bibliography{name}

\begin{enumerate}
    \item Brito e Silva, N. J., et al. (2022). Third- and fifth-order optical nonlinearities of norbixin. \textit{Results in Optics}, \textbf{6}, 100205.
\end{enumerate}
\end{minipage}
}

\end{multicols}

% Acknowledgements
\coloredsection{vibrantblue!60!white}{Acknowledgements}{}

\vspace{0cm}
 
\coloredsubsection{vibrantblue!10!white}{
\Large I would like to acknowledge the professors, staff, and students of the Nonlinear Optics Research and Bio-Optics Research Laboratory who provided support during this work.}


% \textit{Finally, I am very thankful to my wonderful friend \textbf{ Md. Fazle Rabbi Khan} who was beside me all the time, encouraging and inspiring me.}
     


\coloredsubsection{red!50!white}{
\begin{minipage}[t]{\linewidth}
    % \colorbox{highlightgreen!10!white}{%
        % \begin{minipage}{\textwidth}
            \begin{center}{\Large
               \textit{National Conference on Physics for the 21st Century}\\[1ex]
                 \textit{May 18, 2024 } | \textit{Physics Discipline, University name.}}
            \end{center}
\end{minipage}
}

\end{document}


\end{document}